\section{Animation: Additive and Backend-Conditional}
\label{sec:animation}

Animation is modeled as an \emph{additive}, \emph{declarative}, \emph{backend-conditional} layer on top of the static Scene IR. Backends that support animation (SVG, HTML) honor animation hints; backends that produce static output (PNG, PDF, PPTX) render the resting frame.

\paragraph{Implementation Status.}
Fully implemented with optional \texttt{animation} field in Scene IR. SVG backend emits SMIL \texttt{<animate>} and \texttt{<animateMotion>} elements; raster backends (PNG, Skia) render the byte-identical resting frame. Flow grammar edges support flowing dashes (\texttt{stroke-dashoffset} animation). All animation attributes deterministic and included in SVG hash verification.

\subsection{Design Philosophy}

\paragraph{Additive, not fundamental.}
Animation is not required for a diagram to be valid or useful. The vast majority of use cases (print, static slides, documentation) require no animation. Animation is an optional enhancement for web contexts.

\paragraph{Declarative, not frame-based.}
Animation hints describe \emph{what} should animate, not \emph{how} to render each frame. The implementation uses declarative mechanisms (SVG SMIL, CSS animations) that the browser or renderer interprets. There is no frame-by-frame generation, no video encoding, no GIF pipelines.

\paragraph{Backend-conditional.}
Animation hints are carried in the Scene IR. Backends that support animation render them; backends that don't, ignore them. This mirrors the existing pattern for effects (glow is Skia-only; SVG approximates or omits).

\paragraph{Determinism preserved.}
An animated SVG is still byte-identical markup. The animation is expressed as declarative attributes (\texttt{<animate>}, \texttt{<animateMotion>}, CSS \texttt{@keyframes}), not as different frames. Golden tests compare SVG text, including animation elements.

\subsection{Animation Patterns from the Corpus}

Analysis of the inspiration corpus reveals several animation patterns used in technical explainer infographics:

\begin{description}
  \item[Flowing dashes / marching ants] Arrows or connectors with animated \texttt{stroke-dashoffset}, creating a sense of flow or data movement. Common in flow diagrams and architecture visualizations.
  
  \item[Draw-on] Path strokes that animate from zero length to full length, as if being drawn by hand. Used for emphasis when diagrams appear.
  
  \item[Dot along path] A small circle that travels along a connector path, representing data or signal flow.
  
  \item[Pulse / glow cycle] Elements that pulse (opacity or scale) to draw attention. Used sparingly for emphasis.
  
  \item[Sequential reveal] Elements that appear in sequence (fade-in with delay) rather than all at once. Used for step-by-step explanations.
\end{description}

\subsection{Animation Hint Types}

Animation hints are attached to Scene IR primitives via \texttt{animation\_ref} fields:

\subsubsection{FlowingDashes}

Animates \texttt{stroke-dashoffset} to create flowing/marching effect on lines and paths.

\begin{Verbatim}[frame=single,fontsize=\small]
FlowingDashes {
  target_id: string,              // Line or Path id
  dash_length: number,            // px
  gap_length: number,             // px
  speed_px_per_sec: number,       // flow speed
  direction: 'forward' | 'reverse'
}
\end{Verbatim}

\paragraph{SVG implementation.}
\begin{lstlisting}[language=html,caption={SVG flowing dashes animation}]
<line id="flow-1" stroke-dasharray="10 5">
  <animate attributeName="stroke-dashoffset"
           from="0" to="15" dur="1s"
           repeatCount="indefinite"/>
</line>
\end{lstlisting}

\subsubsection{DrawOn}

Animates a path from fully hidden (\texttt{stroke-dashoffset} = path length) to fully visible.

\begin{Verbatim}[frame=single,fontsize=\small]
DrawOn {
  target_id: string,              // Path id
  duration_ms: number,
  easing: 'linear' | 'ease-in-out' | 'ease-out',
  delay_ms?: number
}
\end{Verbatim}

\subsubsection{DotAlongPath}

Animates a dot (circle) traveling along a path using \texttt{<animateMotion>}.

\begin{Verbatim}[frame=single,fontsize=\small]
DotAlongPath {
  target_id: string,              // Path id to follow
  dot_radius: number,
  dot_fill: string,               // color
  duration_ms: number,
  repeat: number | 'infinite'
}
\end{Verbatim}

\subsubsection{Pulse}

Animates opacity, scale, or fill color in a repeating cycle.

\begin{Verbatim}[frame=single,fontsize=\small]
Pulse {
  target_id: string,
  property: 'opacity' | 'scale' | 'fill',
  from: number | string,
  to: number | string,
  duration_ms: number,
  repeat: number | 'infinite'
}
\end{Verbatim}

\subsubsection{FadeIn}

Animates opacity from 0 to 1, optionally with delay for sequential reveal.

\begin{Verbatim}[frame=single,fontsize=\small]
FadeIn {
  target_id: string,
  duration_ms: number,
  delay_ms?: number,
  easing: 'linear' | 'ease-in-out'
}
\end{Verbatim}

\subsection{Theme-Level Animation Defaults}

Themes may specify default animation behavior:

\begin{lstlisting}[language=yaml,caption={Theme animation configuration}]
animation:
  enabled: true                    # master switch
  default_duration_ms: 500
  default_easing: ease-in-out
  connectors:
    style: flowing-dashes          # none | flowing-dashes | draw-on
    speed_px_per_sec: 30
  appear:
    style: fade-in                 # none | fade-in | draw-on
    stagger_ms: 100                # delay between elements
\end{lstlisting}

\subsection{Backend Behavior}

\begin{center}
\small
\begin{tabular}{lll}
\toprule
\textbf{Backend} & \textbf{Animation Support} & \textbf{Behavior} \\
\midrule
SVG & Full & Emits SMIL \texttt{<animate>} / \texttt{<animateMotion>} \\
HTML & Full & Emits CSS \texttt{@keyframes} and \texttt{animation} properties \\
PNG (resvg) & None & Renders resting frame (t=0 or $t=\infty$) \\
PNG (Skia) & None & Renders resting frame \\
PDF & None & Renders resting frame \\
PPTX & Partial & Ignored (PPTX animation is a separate system) \\
\bottomrule
\end{tabular}
\end{center}

\paragraph{Resting frame definition.}
When animation is ignored, the resting frame is:
\begin{itemize}
  \item For \texttt{FlowingDashes}: solid stroke (no dash animation)
  \item For \texttt{DrawOn}: fully visible stroke
  \item For \texttt{DotAlongPath}: dot at start of path
  \item For \texttt{Pulse}: element at \texttt{from} value
  \item For \texttt{FadeIn}: fully opaque ($t=\infty$ state)
\end{itemize}

\subsection{Scope Discipline: What's Out}

The following animation capabilities are explicitly \textbf{out of scope} for v1:

\begin{description}
  \item[Frame-by-frame generation] No rendering of individual animation frames. No GIF/APNG/WebP animated image output.
  
  \item[Video encoding] No MP4/WebM output. Video requires frame rendering, audio timing, and codec complexity.
  
  \item[Lottie/Bodymovin export] Lottie is a rich animation format that would require a separate export path. Deferred to future.
  
  \item[Interactive animation] No animation triggered by user interaction (hover, click). The diagrams are declarative, not interactive.
  
  \item[Complex sequencing] No timeline-based animation sequencing (keyframes at specific times). Only simple, repeating animations.
\end{description}

\paragraph{Rationale.}
These exclusions preserve the elegance of the small-declarative-spec → crisp-vector pipeline. Frame rendering, video encoding, and Lottie export would require heavy dependencies, break the text-based-determinism property, and add significant complexity. They may be valuable future extensions, but they are not v1.
